%% appC --- The tool matrix
%%
%% Written September 2026. Every version here prints from a pinned-version
%% macro in preamble.tex and none is typed: this page is the one place a
%% reader sees every pin at once, and a pin re-typed here would be a second
%% copy of a fact the preamble already owns.
%% tools/check_structure.py --tools fails when a pin macro exists and this
%% appendix does not print it, and when a version is typed rather than
%% printed from a macro.

\chapter{The tool matrix}
\label{app:tools}

\noindent
The version column is not typed. Every number on this page is the same
macro the title page and the chapters print, which is the same string
\code{code/pyproject.toml} pins and \code{uv.lock} resolves \dash{} so the
matrix cannot disagree with what CI installed, and a version bump moves
this page without anybody editing it. A gate fails the build if a pin
exists and this table does not show it.

Read the first column loosely. Some of these are exact counterparts and
some are only the nearest thing a .NET engineer already knows; where the
two differ in a way that matters, the chapter named on the right is where
the difference is spent.

\section{The language and its toolchain}
\label{sec:appC-toolchain}

\begin{toolmatrix}{.NET}{Python}{Verified at}{Ch.}
the runtime & CPython & \pypatch & \ref{ch:runtime} \\
\code{dotnet} CLI, NuGet & \pkg{uv} & \uvver & \ref{ch:packaging} \\
\code{dotnet format}, analysers & \pkg{ruff} & \ruffver &
  \ref{ch:packaging} \\
the compiler's type checking & \pkg{pyright}, strict &
  \pyrightver & \ref{ch:typing} \\
--- & \pkg{mypy}, a second opinion & \mypyver & \ref{ch:typing} \\
\end{toolmatrix}

\noindent
Two checkers are pinned rather than one, and Chapter~\ref{ch:typing}
measures why: on the same four functions pyright reports nothing and mypy
reports two, because one applies the type system and the other adds an
opinion the type system does not have.

\section{Shipping}
\label{sec:appC-shipping}

\begin{toolmatrix}{.NET}{Python}{Verified at}{Ch.}
model binding, validation & \pkg{pydantic} & \pydanticver &
  \ref{ch:typing} \\
\code{IOptions<T>} & \pkg{pydantic-settings} & \pysettingsver &
  \ref{ch:services} \\
ASP.NET Core & \pkg{fastapi} & \fastapiver & \ref{ch:services} \\
Kestrel & \pkg{uvicorn} & \uvicornver & \ref{ch:services} \\
\code{HttpClient} & \pkg{httpx} & \httpxver & \ref{ch:asyncio} \\
Entity Framework & \pkg{sqlalchemy} & \sqlalchemyver & \ref{ch:data} \\
EF migrations & \pkg{alembic} & \alembicver & \ref{ch:data} \\
--- & \pkg{aiosqlite}, for the async listings & \aiosqlitever &
  \ref{ch:data} \\
a DataTable & \pkg{polars} & \polarsver & \ref{ch:data} \\
\end{toolmatrix}

\section{Testing and operations}
\label{sec:appC-testing}

\begin{toolmatrix}{.NET}{Python}{Verified at}{Ch.}
xUnit & \pkg{pytest} & \pytestver & \ref{ch:testing} \\
async test support & \pkg{pytest-asyncio} & \pytestasyncver &
  \ref{ch:testing} \\
FsCheck & \pkg{hypothesis} & \hypothesisver & \ref{ch:testing} \\
coverlet & \pkg{coverage} & \coveragever & \ref{ch:testing} \\
\code{ILogger<T>} & \pkg{structlog} & \structlogver &
  \ref{ch:observability} \\
OpenTelemetry & \pkg{opentelemetry-sdk} & \otelver &
  \ref{ch:observability} \\
\code{dotnet list package -{}-vulnerable} & \pkg{pip-audit} &
  \pipauditver & \ref{ch:observability} \\
\end{toolmatrix}

\section{The AI kit}
\label{sec:appC-ai}

\begin{toolmatrix}{.NET}{Python}{Verified at}{Ch.}
an SDK & \pkg{anthropic} & \anthropicver & \ref{ch:aitoolkit} \\
an SDK & \pkg{openai} & \openaiver & \ref{ch:aitoolkit} \\
\end{toolmatrix}

\begin{note}
Both SDKs depend on the \pkg{httpx2} distribution rather than the
\pkg{httpx} pinned above. They are different packages that share an
ancestry and a spelling, which is why an \api{isinstance} check against the
one you pinned fails on a client that is obviously an httpx client.
Chapter~\ref{ch:aitoolkit} measures it.
\end{note}

\section{Building the book}
\label{sec:appC-book}

\noindent
Two pins are not in the tables above because they are not libraries the
reader installs. The interpreter the book targets is Python~\pyver, and the
exact build every listing ran on is \pypatch, which is what
\code{code/.python-version} names. And the book's own C\# \dash{} the
comparisons its chapters make, and the solutions in
Appendix~\ref{app:interview} \dash{} is compiled and tested in CI against
.NET~\dotnetver, pinned once in \code{code/csharp/global.json} and read
from there by every workflow that needs it.
